% ============================================================
\section{Obtención y procesamiento de los datos}
\label{sec:scripts}
% ============================================================

Para obtener los datos se consultó el índice de ESGF y se identificaron los 102 modelos que publican \texttt{tos}/\texttt{Omon}; se retuvieron los que publican los cuatro experimentos y, de cada uno, la variante de grilla más gruesa y un miembro de ensamble común a los cuatro. Los archivos se descargaron en cascada desde el nodo principal de ESGF, los nodos alternativos de la federación y, como último recurso, Copernicus CDS, y ERSSTv5 se descargó de NOAA PSL. Todos los modelos se regrillaron a la grilla de ERSSTv5, se homogeneizaron calendario y unidades (comprobando el valor numérico y no el atributo declarado), se enmascararon con la máscara océano-tierra de ERSSTv5 y se filtraron valores físicamente imposibles. Finalmente, un control de calidad verificó el rango físico de la TSM ($-2$ a 45\,°C; el límite superior considera que la franja tropical incluye mares marginales cálidos, como el mar Rojo y el golfo Pérsico, donde algunos modelos alcanzan de forma plausible entre 39 y 41\,°C hacia 2100 en SSP3-7.0 y SSP5-8.5) y la cobertura temporal de cada archivo, y solo se seleccionaron los modelos con los cuatro experimentos aprobados. La Figura~\ref{fig:flujograma} resume la secuencia. La descripción de cada paso y de sus controles se encuentra en el documento de anexos.

\begin{figure}[htbp]
\centering
\begin{tikzpicture}[
  scale=0.95, transform shape,
  box/.style={draw, rounded corners, fill=green!80!white!10, align=center, text width=4.2cm, minimum height=.8cm, font=\small},
  boxaux/.style={draw, rounded corners, dashed, fill=green!20!cyan!15, align=center, text width=2.6cm, minimum height=1.1cm, font=\small},
  edg/.style={-{Latex[length=3mm]}, thick},
  edgd/.style={-{Latex[length=3mm]}, dashed, gray!80!black},
  every node/.style={font=\small}
]
\node[box] (S00) at (0,0) {00: Entorno};
\node[box] (S00b) at (0,-1.5) {00b: Modelos candidatos};
\node[boxaux] (S00c) at (5.5,-1.5) {00c\\Verif. literatura};
\node[box] (S01) at (0,-3) {01: Catálogo ESGF};
\node[box] (S02) at (-2.6,-4.5) {02: Descarga CMIP6};
\node[box] (S03) at (3.9,-4.5) {03: Descarga ERSSTv5};
\node[box] (S02b) at (-2.6,-6) {02b: Nodos ESGF alt.};
\node[box] (S02c) at (-2.6,-7.5) {02c: Copernicus CDS};
\node[box] (S04) at (0.5,-9) {04: Grilla común};
\node[box] (S05) at (0.5,-10.5) {05: Máscara océano-tierra};
\node[box] (S06) at (0.5,-12) {06: Control de calidad};
\node[box] (S07) at (0.5,-13.5) {07: Inventario final};
\node[box] (S08) at (0.5,-15) {Registro de modelos};
\node[box] (S09) at (0.5,-16.5) {Gráficos, reporte y paquete};
\draw[edg] (S00) -- (S00b);
\draw[edgd] (S00b) -- (S00c);
\draw[edg] (S00b) -- (S01);
\draw[edg] (S01) -- (S02);
\draw[edg] (S02) -- (S02b);
\draw[edg] (S02b) -- (S02c);
\draw[edgd] (S02b) to[bend right=60] (S01);
\draw[edgd] (S02c) to[bend right=80] (S01);
\draw[edg] (S02.west) -- ++(-1.5,0) |- (S04.west);
\draw[edg] (S03) -- (S04);
\draw[edg] (S04) -- (S05);
\draw[edg] (S05) -- (S06);
\draw[edg] (S06) -- (S07);
\draw[edg] (S07) -- (S08);
\draw[edg] (S08) -- (S09);
\end{tikzpicture}
\caption[Flujograma del \emph{pipeline} del Entregable~1.]{Flujograma del \emph{pipeline}. Las flechas sólidas siguen la secuencia automática de \texttt{run.sh}: buscar los modelos (00b, 01), descargar (02, 03), homogeneizar (04, 05) y verificar (06, 07). Las punteadas indican el diagnóstico manual 00c y la actualización del catálogo cuando un modelo se obtiene de una fuente alternativa (02b, 02c). Las dos últimas cajas se ejecutan en cada corrida.}
\label{fig:flujograma}
\end{figure}
